\pdfoutput=1
\documentclass[11pt]{article}
\usepackage[margin=1in]{geometry}
\usepackage{graphicx,booktabs,amsmath,amssymb,hyperref,natbib,longtable}
\usepackage{cmap}
\usepackage[T1]{fontenc}
\input{glyphtounicode}
\pdfgentounicode=1
\graphicspath{{./}{../results/}}
\hypersetup{colorlinks=true,linkcolor=blue,citecolor=blue,urlcolor=blue,
  pdftitle={Supplementary Material: Chemistry-resolved disagreement between Materials Project-compatible PBE(+U) and r2SCAN stability verdicts in the GNoME materials database},
  pdfauthor={Karthikeya Aduru}}
\setcounter{topnumber}{3}
\setcounter{bottomnumber}{2}
\setcounter{totalnumber}{4}
\renewcommand{\topfraction}{0.9}
\renewcommand{\bottomfraction}{0.7}
\renewcommand{\textfraction}{0.07}
\renewcommand{\floatpagefraction}{0.85}
\title{\textbf{Supplementary Material for:\\ Chemistry-resolved disagreement between Materials Project-compatible
PBE(+$U$) and r$^2$SCAN stability verdicts in the GNoME materials database}}
\author{Karthikeya Aduru \\ \small Accretion Labs \\ \small \texttt{founder@accretionlabs.net}}
\date{October 2026}

\renewcommand{\thesection}{S\arabic{section}}
\renewcommand{\thetable}{S\arabic{table}}
\begin{document}
\maketitle
\noindent This supplement accompanies the main text and holds the cell-level, per-element
and matched-cell tables and the input manifest. Section and table numbers prefixed
``main-text'' refer to the main paper; every printed value here is checked against the
analysis output by the same \texttt{check\_paper.py} that checks the main text.
\sloppy

\section{Cell-level results with cluster bootstrap intervals}\label{app:cells}

This section lists all disjoint cells with $n \ge 500$. Intervals resample prototype
clusters. The final column is the ratio of the cluster interval width to the naive Wald binomial width.

\footnotesize
\begin{longtable}{llccrrrr}
\caption{Flip rate by composition cell, all disjoint cells with $n \ge 500$, with cluster bootstrap 95\% intervals and the widening factor relative to the naive binomial interval.}\label{tab:appcells}\\
\toprule
$n_{\text{el}}$ & Anion & Ln & 3$d$ & $n$ & Rate & Cluster 95\% CI & Widen \\
\midrule
\endfirsthead
\toprule
$n_{\text{el}}$ & Anion & Ln & 3$d$ & $n$ & Rate & Cluster 95\% CI & Widen \\
\midrule
\endhead
4 & oxide & no & yes & 2{,}478 & 30.7\% & [27.7, 33.8] & 1.69 \\
3 & chalcogenide & no & yes & 564 & 30.3\% & [24.9, 36.5] & 1.53 \\
4 & pnictide & yes & yes & 4{,}573 & 30.1\% & [25.9, 33.8] & 2.96 \\
3 & boride/carbide/silicide & no & yes & 575 & 29.6\% & [24.1, 35.3] & 1.51 \\
4 & boride/carbide/silicide & no & yes & 4{,}501 & 25.2\% & [21.5, 29.5] & 3.14 \\
3 & pnictide & no & yes & 813 & 24.0\% & [19.9, 28.7] & 1.51 \\
4 & boride/carbide/silicide & yes & yes & 4{,}561 & 22.8\% & [18.5, 27.4] & 3.65 \\
4 & halide & no & yes & 794 & 22.4\% & [17.5, 27.9] & 1.80 \\
3 & oxide & no & no & 921 & 21.7\% & [18.7, 24.9] & 1.17 \\
3 & intermetallic & no & yes & 2{,}667 & 21.1\% & [18.4, 24.0] & 1.80 \\
3 & intermetallic & yes & yes & 2{,}056 & 19.0\% & [15.9, 22.4] & 1.91 \\
4 & oxide & no & no & 3{,}318 & 18.9\% & [16.7, 21.4] & 1.79 \\
4 & oxide & yes & yes & 1{,}399 & 18.1\% & [14.2, 22.0] & 1.94 \\
4 & pnictide & no & yes & 5{,}736 & 17.8\% & [14.7, 21.0] & 3.18 \\
3 & halide & no & no & 1{,}326 & 17.3\% & [14.3, 20.8] & 1.60 \\
3 & pnictide & yes & no & 1{,}349 & 16.8\% & [11.9, 22.4] & 2.63 \\
4 & chalcogenide & no & yes & 1{,}487 & 16.3\% & [13.4, 19.7] & 1.67 \\
4 & pnictide & yes & no & 7{,}412 & 13.9\% & [11.5, 16.5] & 3.15 \\
4 & oxide & yes & no & 4{,}232 & 13.8\% & [11.4, 16.6] & 2.52 \\
4 & halide & no & no & 1{,}404 & 13.1\% & [10.2, 16.6] & 1.82 \\
4 & intermetallic & yes & yes & 10{,}659 & 12.9\% & [11.6, 14.5] & 2.27 \\
3 & oxide & yes & no & 730 & 12.2\% & [7.7, 16.9] & 1.93 \\
4 & intermetallic & no & yes & 7{,}459 & 12.1\% & [10.7, 13.5] & 1.88 \\
3 & boride/carbide/silicide & no & no & 845 & 11.7\% & [9.0, 14.9] & 1.36 \\
3 & pnictide & no & no & 1{,}473 & 11.1\% & [8.9, 13.6] & 1.47 \\
3 & halide & yes & no & 982 & 10.6\% & [6.7, 15.8] & 2.37 \\
3 & chalcogenide & no & no & 1{,}106 & 10.6\% & [8.4, 13.0] & 1.26 \\
3 & chalcogenide & yes & no & 1{,}076 & 10.1\% & [7.4, 13.7] & 1.77 \\
3 & boride/carbide/silicide & yes & no & 1{,}246 & 9.6\% & [6.1, 13.9] & 2.37 \\
4 & halide & yes & yes & 1{,}119 & 9.1\% & [5.3, 14.5] & 2.74 \\
4 & boride/carbide/silicide & yes & no & 10{,}287 & 9.1\% & [7.0, 11.4] & 4.02 \\
4 & chalcogenide & yes & no & 1{,}497 & 8.4\% & [6.3, 10.7] & 1.59 \\
4 & boride/carbide/silicide & no & no & 3{,}422 & 8.2\% & [6.8, 9.7] & 1.59 \\
3 & intermetallic & no & no & 2{,}618 & 7.9\% & [6.6, 9.3] & 1.30 \\
4 & chalcogenide & no & no & 1{,}467 & 7.0\% & [5.5, 8.7] & 1.22 \\
3 & intermetallic & yes & no & 4{,}661 & 6.4\% & [5.0, 8.0] & 2.11 \\
4 & pnictide & no & no & 4{,}675 & 6.0\% & [4.8, 7.3] & 1.82 \\
4 & halide & yes & no & 3{,}599 & 5.9\% & [2.8, 10.8] & 5.16 \\
4 & intermetallic & yes & no & 18{,}932 & 5.5\% & [4.6, 6.3] & 2.59 \\
4 & intermetallic & no & no & 5{,}025 & 4.0\% & [3.3, 4.7] & 1.31 \\
\bottomrule
\end{longtable}
\normalsize

The median widening factor is 1.82. The highest and lowest quaternary cells remain disjoint:
30.7\% [27.7, 33.8] against 4.0\% [3.3, 4.7].

\section{Per-element results}\label{app:element}

This section lists all 83 elements appearing in at least 300 recomputed materials,
ordered by flip rate and read left to right across the three column groups. Entries are point estimates; cluster bootstrap
intervals for the extremes are given in main-text Section~4.2. An asterisk marks the
53 elements whose contrast against the global rate survives Benjamini--Hochberg
correction at 0.05. Each element is contrasted against a global rate that includes that
element, which attenuates the contrast for the largest-$n$ entries.

\footnotesize
\begin{longtable}{lrr@{\hspace{1.6em}}lrr@{\hspace{1.6em}}lrr}
\caption{Per-element flip rate for the 83 elements with $n \ge 300$; an asterisk marks a Benjamini--Hochberg-significant contrast against the global rate.}\label{tab:appelement}\\
\toprule
El. & $n$ & Flip (\%) & El. & $n$ & Flip (\%) & El. & $n$ & Flip (\%) \\
\midrule
\endfirsthead
\toprule
El. & $n$ & Flip (\%) & El. & $n$ & Flip (\%) & El. & $n$ & Flip (\%) \\
\midrule
\endhead
Fe$^{*}$ & 5{,}762 & 42.8 & Co$^{*}$ & 10{,}972 & 33.2 & Cr$^{*}$ & 2{,}553 & 30.7 \\
Mn$^{*}$ & 4{,}994 & 25.2 & F$^{*}$ & 2{,}801 & 24.9 & As$^{*}$ & 6{,}298 & 21.6 \\
Ce$^{*}$ & 4{,}236 & 20.5 & Mo$^{*}$ & 3{,}910 & 20.2 & Sb$^{*}$ & 6{,}577 & 19.9 \\
O$^{*}$ & 13{,}668 & 19.3 & V$^{*}$ & 3{,}659 & 18.9 & Si$^{*}$ & 18{,}382 & 18.4 \\
W$^{*}$ & 3{,}223 & 18.2 & Eu$^{*}$ & 2{,}132 & 18.0 & U$^{*}$ & 2{,}595 & 17.9 \\
Te$^{*}$ & 2{,}207 & 17.4 & Bi$^{*}$ & 6{,}606 & 16.5 & Ag$^{*}$ & 4{,}224 & 16.3 \\
Gd & 976 & 16.0 & Zr & 7{,}315 & 14.7 & P & 10{,}265 & 14.2 \\
Sc & 7{,}842 & 13.9 & Se & 4{,}543 & 13.6 & Ca & 5{,}153 & 13.5 \\
Au & 5{,}932 & 13.5 & S & 3{,}383 & 13.4 & Ni & 10{,}142 & 13.4 \\
Cl & 2{,}172 & 13.2 & Y & 11{,}046 & 13.1 & Ti & 5{,}612 & 13.1 \\
Hf & 6{,}809 & 13.0 & C & 4{,}735 & 12.9 & Ga & 10{,}930 & 12.9 \\
Dy & 11{,}559 & 12.8 & K & 2{,}795 & 12.7 & H & 2{,}416 & 12.7 \\
Cu & 6{,}254 & 12.6 & Ho & 11{,}988 & 12.6 & Pr & 8{,}238 & 12.4 \\
Sr & 3{,}981 & 12.4 & Sn & 9{,}436 & 12.3 & Lu & 6{,}821 & 12.3 \\
Li & 6{,}305 & 12.2 & Ba & 3{,}580 & 12.1 & Tb$^{*}$ & 12{,}248 & 12.1 \\
Nb & 5{,}297 & 11.9 & Tm$^{*}$ & 10{,}154 & 11.8 & Er$^{*}$ & 10{,}434 & 11.6 \\
Na & 2{,}932 & 11.6 & Ta$^{*}$ & 5{,}021 & 11.3 & Nd$^{*}$ & 9{,}238 & 11.1 \\
La$^{*}$ & 6{,}546 & 11.0 & Rb$^{*}$ & 3{,}141 & 10.8 & Pu & 1{,}020 & 10.8 \\
Br & 3{,}440 & 10.6 & Pb$^{*}$ & 4{,}377 & 10.6 & B$^{*}$ & 9{,}496 & 10.6 \\
Hg$^{*}$ & 4{,}445 & 10.2 & Cs$^{*}$ & 2{,}973 & 10.1 & Th$^{*}$ & 7{,}632 & 9.9 \\
Al$^{*}$ & 11{,}933 & 9.7 & Np$^{*}$ & 1{,}387 & 9.7 & Zn$^{*}$ & 7{,}510 & 9.1 \\
Cd$^{*}$ & 5{,}356 & 9.0 & In$^{*}$ & 8{,}259 & 9.0 & Os$^{*}$ & 8{,}652 & 8.9 \\
N$^{*}$ & 3{,}845 & 8.7 & Pd$^{*}$ & 10{,}381 & 8.1 & Re$^{*}$ & 3{,}958 & 8.1 \\
Ru$^{*}$ & 12{,}381 & 8.1 & Sm$^{*}$ & 4{,}852 & 8.0 & I$^{*}$ & 3{,}091 & 7.9 \\
Pt$^{*}$ & 13{,}450 & 7.8 & Ge$^{*}$ & 4{,}764 & 7.6 & Rh$^{*}$ & 14{,}340 & 7.3 \\
Yb$^{*}$ & 6{,}265 & 7.1 & Be$^{*}$ & 2{,}799 & 6.5 & Ir$^{*}$ & 15{,}680 & 6.2 \\
Tc$^{*}$ & 7{,}350 & 5.7 & Mg$^{*}$ & 4{,}073 & 5.7 & Pm$^{*}$ & 389 & 4.4 \\
Pa$^{*}$ & 4{,}251 & 3.6 & Ac$^{*}$ & 2{,}181 & 1.4 &  & &  \\
\bottomrule
\end{longtable}
\normalsize

\section{All matched cells against known materials}\label{app:allcells}

Table~\ref{tab:appallcells} lists all twenty cells of main-text Section~4.11: the ten of
main-text Table~11 together with the ten whose intervals straddle zero, under the
same cell thresholds, bootstrap and conventions as that table. Positive $\delta$ means
GNoME disagrees more; $\ast$ marks the ten that exclude zero. This table, that table
and the standardised contrasts of the same section are all generated by
\texttt{phase37\_matched\_cluster.py}.

\begin{table}[htbp]\centering
\caption{All twenty matched cells, GNoME against Materials Project previously catalogued materials; conventions as in main-text Table~11; $\delta$ is computed from unrounded rates and $\ast$ marks intervals excluding zero.}\label{tab:appallcells}
\small\begin{tabular}{lrrrrrl}
\toprule
Cell & $n_{\text{G}}$ & GNoME & $n_{\text{K}}$ & Known & $\delta$ & 95\% CI \\
\midrule
4-el oxide, 3$d$ & 2{,}478 & 30.7\% & 180 & 13.3\% & $+17.3$ & [$+11.4$, $+23.0$] $\ast$ \\
3-el chalcogenide, 3$d$ & 564 & 30.3\% & 450 & 17.8\% & $+12.5$ & [$+6.0$, $+19.4$] $\ast$ \\
3-el bor./carb./sil., 3$d$ & 575 & 29.6\% & 454 & 18.3\% & $+11.3$ & [$+4.6$, $+18.0$] $\ast$ \\
3-el halide, no Ln/3$d$ & 1{,}326 & 17.3\% & 1{,}125 & 9.2\% & $+8.2$ & [$+4.7$, $+11.9$] $\ast$ \\
4-el oxide, no Ln/3$d$ & 3{,}318 & 18.9\% & 297 & 10.8\% & $+8.2$ & [$+3.9$, $+12.4$] $\ast$ \\
3-el oxide, no Ln/3$d$ & 921 & 21.7\% & 1{,}553 & 16.7\% & $+5.0$ & [$+1.4$, $+8.6$] $\ast$ \\
3-el pnictide, no Ln/3$d$ & 1{,}473 & 11.1\% & 852 & 6.9\% & $+4.1$ & [$+1.3$, $+7.1$] $\ast$ \\
3-el pnictide, 3$d$ & 813 & 24.0\% & 557 & 20.3\% & $+3.7$ & [$-1.7$, $+9.5$] \\
3-el chalcogenide, no Ln/3$d$ & 1{,}106 & 10.6\% & 1{,}064 & 7.4\% & $+3.2$ & [$+0.5$, $+6.0$] $\ast$ \\
3-el bor./carb./sil., Ln & 1{,}246 & 9.6\% & 674 & 7.1\% & $+2.5$ & [$-1.5$, $+7.2$] \\
3-el chalcogenide, Ln & 1{,}076 & 10.1\% & 442 & 7.9\% & $+2.2$ & [$-1.6$, $+6.5$] \\
3-el bor./carb./sil., no Ln/3$d$ & 845 & 11.7\% & 448 & 9.6\% & $+2.1$ & [$-1.8$, $+6.1$] \\
4-el oxide, Ln & 4{,}232 & 13.8\% & 108 & 12.0\% & $+1.8$ & [$-5.2$, $+8.2$] \\
3-el pnictide, Ln & 1{,}349 & 16.8\% & 422 & 15.4\% & $+1.4$ & [$-4.7$, $+8.0$] \\
3-el oxide, Ln & 730 & 12.2\% & 669 & 12.7\% & $-0.5$ & [$-5.8$, $+4.8$] \\
3-el intermetallic, Ln, 3$d$ & 2{,}056 & 19.0\% & 1{,}153 & 20.6\% & $-1.5$ & [$-5.4$, $+2.5$] \\
3-el intermetallic, 3$d$ & 2{,}667 & 21.1\% & 1{,}125 & 23.1\% & $-2.0$ & [$-5.7$, $+1.8$] \\
3-el halide, Ln & 982 & 10.6\% & 434 & 14.8\% & $-4.2$ & [$-9.3$, $+1.9$] \\
3-el intermetallic, Ln & 4{,}661 & 6.4\% & 1{,}988 & 13.0\% & $-6.6$ & [$-8.7$, $-4.5$] $\ast$ \\
3-el intermetallic, no Ln/3$d$ & 2{,}618 & 7.9\% & 1{,}144 & 15.7\% & $-7.7$ & [$-10.2$, $-5.3$] $\ast$ \\
\bottomrule
\end{tabular}
\end{table}

\section{Input manifest}\label{app:manifest}

\paragraph{Licensing.} GNoME data are released under the Creative Commons
Attribution-NonCommercial 4.0 International licence and the accompanying code under
Apache 2.0 \citep{gnomedataset-s}; they are used here for non-commercial research. This
paper's repository redistributes no GNoME records: it holds code, input checksums and
derived aggregate statistics, the largest GNoME-derived result file carrying 83 rows of
per-element summaries. The two per-material files it does commit are Materials Project
extracts, used under that project's Creative Commons Attribution 4.0 terms with the
attribution recorded in \texttt{results/PROVENANCE.md} \citep{jain2013-s}.

\paragraph{Files.} Two CSV files were downloaded from
\texttt{gs://gdm\_materials\_discovery/gnome\_data/} on 9 October 2026; the checksums
are of the files as downloaded, before any processing.

\noindent\texttt{stable\_materials\_summary.csv}: 554{,}054 rows, 151{,}452{,}777 bytes, SHA-256\\
\texttt{c37bffb9dd6c51d38d259a2d340ecf796194729f656d474828512390db1ff0c9}

\noindent\texttt{stable\_materials\_r2scan.csv}: 139{,}541 rows, 22{,}004{,}293 bytes, SHA-256\\
\texttt{6a811395b065f43ae7911597405eec0ecfa88251e918f4c0ba8e4e99a75656fd}

\paragraph{Columns, missing values and joins.} From the summary file we read
\texttt{MaterialId}, \texttt{Decomposition Energy Per Atom MP}, \texttt{Decomposition
Energy Per Atom All}, \texttt{Elements}, \texttt{Composition}, \texttt{NSites},
\texttt{Bandgap}, \texttt{Crystal System}, \texttt{Data Directory} and \texttt{Is Train}.
From the r$^2$SCAN file we read the same identifier, element and composition fields,
its \texttt{Decomposition Energy Per Atom MP}, \texttt{MP OQMD} and \texttt{All}
columns, and \texttt{Space Group Number}, \texttt{Crystal System} and \texttt{NSites}. Decomposition
energies are in eV/atom and band gaps in eV, as released. \texttt{Elements} is parsed
with \texttt{ast.literal\_eval} into a set; \texttt{Composition} is parsed with the
regular expression \texttt{([A-Z][a-z]?)(\textbackslash d+)} and is \emph{not} reduced.
Energy and gap fields are coerced with \texttt{pandas.to\_numeric(errors="coerce")};
rows with a missing r$^2$SCAN decomposition energy are dropped before any analysis, and
the band-gap accounting of main-text Section~4.7 separates missing, infinite, zero and
positive values. The two files are joined on \texttt{MaterialId}, which is unique in both
files, so no duplicate resolution is required. The 22{,}498 r$^2$SCAN identifiers absent
from the summary file are consistent with a release-version mismatch, which we state as
a hypothesis not confirmed with the dataset authors.

\paragraph{Materials Project queries.} Thermo records were retrieved through the public
API on 9--10 October 2026 from
\url{https://api.materialsproject.org/materials/thermo/}, paging at
\texttt{\_per\_page=1000}. The \texttt{thermo\_types} parameter is passed once per pull
with the literal values \texttt{r2SCAN} and \texttt{GGA\_GGA+U} exactly as spelled; the
capitalisation is the API's and is not interchangeable, and we observed the upper-case
\texttt{R2SCAN} of the client's \texttt{ThermoType} enumeration return zero documents
without raising an error, so the spelling is worth pinning in any replication. The
\texttt{\_fields} parameter is the literal string
\texttt{material\_id,thermo\_type,energy\_above\_hull,nelements,chemsys}; no other filter
is applied. Stability is defined as $E_{\text{hull}} \le 10^{-6}$ eV/atom on each side.
Records are collapsed to one value per \texttt{material\_id}; we checked the cached
responses and found zero duplicate \texttt{(material\_id, thermo\_type)} combinations,
152{,}426 unique identifiers of 152{,}426 rows for \texttt{GGA\_GGA+U} and 37{,}123 of
37{,}123 for r$^2$SCAN. The pulls are cached in \texttt{results/mp\_gga\_raw.csv} and
\texttt{results/mp\_r2scan\_raw.csv}; they are the only components of the pipeline that
require credentials, and the stored copies are reused by the continuous-integration run.


\begin{thebibliography}{9}

\bibitem[Google DeepMind(2023)]{gnomedataset-s} Google DeepMind, ``GNoME dataset
documentation'' (2023), \texttt{DATASET.md} and \texttt{README.md}, \url{https://github.com/google-deepmind/materials_discovery},
accessed 11 October 2026. Data released under CC BY-NC 4.0.

\bibitem[Jain et al.(2013)]{jain2013-s} A.~Jain, S.~P. Ong, G.~Hautier, W.~Chen,
W.~D. Richards, S.~Dacek, S.~Cholia, D.~Gunter, D.~Skinner, G.~Ceder and K.~A. Persson,
``Commentary: The Materials Project: A materials genome approach to accelerating materials
innovation,'' \emph{APL Mater.} \textbf{1}, 011002 (2013). doi:10.1063/1.4812323

\end{thebibliography}
\end{document}
